\documentclass[10pt]{article}
\providecommand{\FIGDIR}{fig}
\newcommand{\DIGESTDATE}{13 August 2026}
% ===================== PM Research Watch - shared preamble =====================
\usepackage[T1]{fontenc}
\usepackage[utf8]{inputenc}
\usepackage{amsmath,amssymb}
\usepackage{tgpagella}
\usepackage{tgheros}
\renewcommand{\sfdefault}{qhv}
\renewcommand{\ttdefault}{lmtt}
\usepackage{microtype}
\usepackage[english]{babel}

\usepackage[a4paper,top=17mm,bottom=17mm,left=17mm,right=17mm,
            headsep=6mm,footskip=10mm]{geometry}
\usepackage{graphicx}
\usepackage{xcolor}
\usepackage{colortbl}
\usepackage{tikz}
\usetikzlibrary{calc}
\usepackage[most]{tcolorbox}
\usepackage{titlesec}
\usepackage{fancyhdr}
\usepackage{booktabs}
\usepackage{array}
\usepackage{longtable}
\usepackage{enumitem}
\usepackage{multicol}
\usepackage{ragged2e}
\usepackage[hidelinks]{hyperref}

% ------------------------------------------------------------------ palette
\definecolor{Ink}    {HTML}{14262E}
\definecolor{Deep}   {HTML}{0F4C5C}
\definecolor{Teal}   {HTML}{1F7A8C}
\definecolor{Sky}    {HTML}{4FA3C4}
\definecolor{Sage}   {HTML}{6FA287}
\definecolor{Amber}  {HTML}{E3A008}
\definecolor{Clay}   {HTML}{B5643C}
\definecolor{Coral}  {HTML}{D1495B}
\definecolor{Violet} {HTML}{6A5B9E}
\definecolor{Slate}  {HTML}{4A6572}
\definecolor{Paper}  {HTML}{FBF9F5}
\definecolor{Mist}   {HTML}{EEEAE2}
\definecolor{Rule}   {HTML}{D9D3C9}

\pagecolor{Paper}
\color{Ink}

\hypersetup{colorlinks=true, linkcolor=Deep, urlcolor=Teal, citecolor=Teal,
            pdftitle={PM Research Watch}, pdfauthor={Automated daily digest}}

% ------------------------------------------------------------------ headings
\titleformat{\section}
  {\sffamily\bfseries\large\color{Deep}}{}{0pt}
  {}[{\vspace{1.2mm}\color{Amber}\titlerule[1.1pt]}]
\titlespacing*{\section}{0pt}{5.5mm}{2.8mm}

\titleformat{\subsection}
  {\sffamily\bfseries\normalsize\color{Teal}}{}{0pt}{}
\titlespacing*{\subsection}{0pt}{3.6mm}{1.4mm}

\setlist[itemize]{leftmargin=4.2mm, itemsep=1.1mm, topsep=1.3mm, parsep=0pt,
                  label=\textcolor{Teal}{\footnotesize$\blacktriangleright$}}

% ------------------------------------------------------------------ header/footer
\pagestyle{fancy}
\fancyhf{}
\renewcommand{\headrulewidth}{0pt}
\renewcommand{\footrulewidth}{0pt}
\fancyhead[L]{\sffamily\scriptsize\color{Slate}PM RESEARCH WATCH\;\textcolor{Rule}{\textbar}\;\DIGESTDATE}
\fancyhead[R]{\sffamily\scriptsize\color{Slate}Particulate matter monitoring \& health effects}
\fancyfoot[C]{\sffamily\scriptsize\color{Slate}\thepage}

% ------------------------------------------------------------------ components
% metric tile
\newcommand{\metric}[3]{%
  \begin{tcolorbox}[enhanced, colback=white, colframe=Rule, boxrule=0.5pt,
    arc=1.6mm, left=2.2mm, right=2.2mm, top=1.8mm, bottom=1.8mm,
    borderline west={1.7mm}{0pt}{#3}, nobeforeafter, valign=center,
    height=17mm, width=\linewidth]
    \RaggedRight
    {\sffamily\bfseries\fontsize{17}{17}\selectfont\color{#3}#1}\\[0.5mm]
    {\sffamily\fontsize{6.5}{7.4}\selectfont\color{Slate}#2}
  \end{tcolorbox}}

% pull-quote / key insight box
\newtcolorbox{keybox}[1][]{enhanced, colback=Mist, colframe=Mist,
  arc=1.6mm, left=3.2mm, right=3.2mm, top=2.4mm, bottom=2.4mm,
  borderline west={1.7mm}{0pt}{Amber}, #1}

% subtopic band
\newcommand{\band}[2]{%
  \vspace{3.2mm}\noindent
  \begin{tcolorbox}[enhanced, colback=#2, colframe=#2, arc=1.2mm,
    left=2.6mm, right=2.6mm, top=1.1mm, bottom=1.1mm, boxrule=0pt,
    width=\linewidth, nobeforeafter]
    {\sffamily\bfseries\small\color{white}#1}
  \end{tcolorbox}\vspace{1.6mm}}

% one paper entry: {short cite}{title}{journal}{doi}{body}
\newcommand{\paperentry}[5]{%
  \par\vspace{1.9mm}\noindent
  {\sffamily\bfseries\footnotesize\color{Ink}#1}\;%
  {\sffamily\fontsize{7}{8}\selectfont\color{Slate}#3}\\[0.4mm]
  {\footnotesize\itshape\color{Deep}#2}\\[0.6mm]
  {\footnotesize #4}\\[0.3mm]
  {\fontsize{6.6}{7.6}\selectfont\color{Slate}\texttt{doi:}\,\href{https://doi.org/#5}{\texttt{\detokenize{#5}}}}
  \par}

% caption style
\newcommand{\figcap}[1]{\par\vspace{1.0mm}
  {\fontsize{7.2}{8.6}\selectfont\color{Slate}\textsf{\textbf{Fig.}}\; #1\par}\vspace{1.5mm}}

\providecommand{\permil}{\textperthousand}

\setlength{\parindent}{0pt}
\setlength{\parskip}{0pt}
\renewcommand{\arraystretch}{1.18}


\begin{document}

% ==================================================================== MASTHEAD
\thispagestyle{empty}
\begin{tikzpicture}[remember picture, overlay]
  \fill[Deep] (current page.north west) rectangle
      ($(current page.north east)+(0,-67mm)$);
  \fill[Amber] ($(current page.north west)+(0,-67mm)$) rectangle
      ($(current page.north east)+(0,-68.4mm)$);
  \foreach \x/\y/\r in {22/12/0.9, 41/26/1.5, 63/9/0.7, 88/31/1.1, 112/17/1.9,
                        138/29/0.8, 159/13/1.3, 178/54/1.0, 196/21/1.6,
                        30/44/0.6, 74/57/1.2, 125/48/0.9, 168/6/0.8, 191/59/0.7,
                        52/18/1.0, 100/60/0.6, 148/20/1.1, 66/49/1.4, 205/45/0.8}
    \fill[Sky, opacity=0.30] ($(current page.north west)+(\x mm,-\y mm)$) circle (\r mm);
\end{tikzpicture}

\vspace*{4mm}
{\sffamily\fontsize{7.6}{9}\selectfont\color{Sky}%
 AUTOMATED DAILY LITERATURE DIGEST\;\textbullet\;PhD RESEARCH WATCH\par}
\vspace{2.2mm}
{\sffamily\bfseries\fontsize{27}{29}\selectfont\color{white}%
 Particulate Matter\par}
{\sffamily\fontsize{27}{29}\selectfont\color{Sky}%
 Monitoring \& Health Effects\par}
\vspace{5.0mm}
{\sffamily\fontsize{8.4}{10}\selectfont\color{white}%
 Issue 2026-08-13\;\;\textcolor{Amber}{\textbar}\;\;Entry window 13 August 2026\;\;%
 \textcolor{Amber}{\textbar}\;\;17 records carried, 12 rejected\par}
\vspace{18mm}

% -------------------------------------------------------------- metric strip
\noindent
\begin{minipage}[t]{0.190\linewidth}\metric{17}{RECORDS CARRIED\\(12 REJECTED)}{Deep}\end{minipage}\hfill
\begin{minipage}[t]{0.190\linewidth}\metric{8}{SUBTOPIC CLUSTERS\\REPRESENTED}{Teal}\end{minipage}\hfill
\begin{minipage}[t]{0.190\linewidth}\metric{4}{EFFECT ESTIMATES\\WITH INTERVALS}{Sky}\end{minipage}\hfill
\begin{minipage}[t]{0.190\linewidth}\metric{9}{INSTRUMENTATION /\\EXPOSURE RECORDS}{Amber}\end{minipage}\hfill
\begin{minipage}[t]{0.190\linewidth}\metric{0}{METADATA-ONLY\\(NO ABSTRACT)}{Clay}\end{minipage}

\vspace{6mm}

% -------------------------------------------------------------- signal of day
\section*{\sffamily\bfseries\color{Deep}Signal of the day}
\vspace{-1.5mm}{\color{Amber}\rule{\linewidth}{1.1pt}}\vspace{2.5mm}

\begin{keybox}
{\sffamily\bfseries\small\color{Deep}A low-cost sensor calibration is only valid inside the
box it was trained in --- and the failure mode differs by which wall you cross.}\\[1.4mm]
\footnotesize
Benjamin et al.\ (\textit{Sensors}, Hannigan group at Colorado) ran a multi-season
co-location of a Plantower PMS5003 and an Alphasense CO-B4 against Federal Equivalence
Method monitors, then deliberately trained calibration models on \emph{restricted} ranges
of relative humidity, temperature and pollutant concentration and scored them on data
outside those ranges. Three results, and the third is the one that is not obvious:
within the training envelope, bias error was minimal; extrapolating outside it
introduced \textbf{bias} error; and extrapolating specifically into \emph{higher
concentrations} than the training data increased \textbf{random} error.

That asymmetry matters operationally. Bias is in principle correctable after the fact if
the excursion can be characterised; inflated variance at high concentration is not, and
high concentration is exactly where a network is asked to support an exceedance call or a
smoke advisory. The design rule follows: a training set must span the \emph{upper} tail of
the expected concentration distribution even more than it spans the humidity range,
because the humidity penalty appears as a correctable offset and the concentration penalty
appears as noise. Transferability claims for LCS calibration need stating as an envelope,
not as an $R^2$.
\end{keybox}

\vspace{3mm}

% -------------------------------------------------------------- exec brief
\section*{\sffamily\bfseries\color{Deep}What changed today}
\vspace{-1.5mm}{\color{Amber}\rule{\linewidth}{1.1pt}}\vspace{2.5mm}

\footnotesize
\begin{itemize}
\item \textbf{An unusually instrumentation-heavy day: 9 of 17 records.} After a fortnight
dominated by burden estimates and reviews, today's harvest is led by sensing and exposure
work --- LCS calibration envelopes, a self-regenerating acoustic particle sensor, a
broadband-sensor-array approach to VOC speciation, gradient-boosted PM$_{2.5}$
forecasting, Pandora column NO$_2$ sub-pixel variability, and a satellite-monitoring
scoping review for Africa.

\item \textbf{Component-resolved PM keeps outperforming mass --- twice, independently.}
Xu et al.\ report cerebral hypoperfusion ORs of \textbf{4.344 (3.323--5.734)} for
PM$_{2.5}$ and \textbf{4.396 (3.327--5.869)} for black carbon in 967 patients with
symptomatic intracranial stenosis, with organic matter and BC together carrying
\textbf{86\% of the WQS mixture effect}. Tan et al.\ find the PM$_{2.5}$ component
mixture associated with lower eGFR in 8{,}477 Wuhan subway workers (WQS OR \textbf{1.183,
1.106--1.266}). Both used the TAP component product; neither could measure components
directly. The measurement gap this exposes is now the recurring theme of this watch.

\item \textbf{Sensor fouling gets an engineering answer.} Fawaz et al.\ actuated a water
droplet with Rayleigh surface acoustic waves to sweep collected particles off the sensing
zone of SAW microparticle sensors inside a cascade impactor, for both NaCl solution and
SiC particles, without degrading sensor performance. Surface saturation is the limiting
failure mode for gravimetric-equivalent sensors in long deployments; a self-regenerating
prototype is a materially different maintenance model.

\item \textbf{Senolytics reverse PM-induced endothelial activation in vivo.} Gomes et
al.\ exposed C57BL/6J mice to concentrated ambient particles and observed increased
leukocyte rolling and adhesion on intact vasculature plus senescence of peripheral
mononuclear and endothelial progenitor cells. Dasatinib + quercetin reversed the
senescence markers ($\beta$-galactosidase, SASP genes) and restored EPC tube-forming
capacity. Mechanistic rather than translational --- but it places senescence upstream of
early atherogenesis rather than downstream of it.

\item \textbf{Africa's satellite-monitoring evidence base is 26 records.} Gould \&
Jagarnath's PRISMA-ScR scoping review found 26 eligible studies linking satellite-derived
air quality to health outcomes across the entire continent, mostly post-2015, mostly
ecological, concentrated in West Africa. First record in this watch to put a number on a
gap its geography figure has shown as an empty bar since 27 July.

\item \textbf{No metadata-only records today} --- first clean issue in five. The 08-08 to
08-12 Elsevier backlog is unchanged and stays queued for Saturday.
\end{itemize}

\clearpage

% ==================================================================== ANALYTICS
\section*{\sffamily\bfseries\color{Deep}Corpus analytics}
\vspace{-1.5mm}{\color{Amber}\rule{\linewidth}{1.1pt}}\vspace{3mm}

\noindent\includegraphics[width=\linewidth]{\FIGDIR/f1_subtopics.png}
\figcap{Subtopic assignment of the 17 carried records, single-label by dominant
contribution. \textit{Sensing, forecasting \& instrumentation} is the largest cluster at
7 records --- the highest single-day count this watch has recorded --- and with
\textit{Exposure assessment} accounts for 9 of 17. \textit{Respiratory \& allergic} and
\textit{Reproductive \& developmental} are both empty today.}

\vspace{2mm}
\noindent\includegraphics[width=\linewidth]{\FIGDIR/f2_design_endpoint.png}
\figcap{Left --- study architecture. Modelling and algorithm work plus instrument
development together dominate; no randomised or interventional design appears.
Right --- health endpoint. Nine records carry no human health endpoint, which is the
direct consequence of the instrumentation share.}

\vspace{2mm}
\noindent\includegraphics[width=\linewidth]{\FIGDIR/f3_metric_geography.png}
\figcap{Left --- particle metric. Note how many of today's exposure metrics are not PM
mass: total column NO$_2$, aerosol depolarization ratio, VOC sensor response, a smoke
plume in an image. On an instrumentation-heavy day the corpus stops being about mass
concentration almost entirely. Right --- geographic provenance. China contributes six
records; a Research Square scoping review is the only African-focused record and it is a
review of how little there is.}

\vspace{2mm}
\noindent\includegraphics[width=\linewidth]{\FIGDIR/f6_lifecourse.png}
\figcap{Life-course windows addressed. Counts are non-exclusive and describe this issue
only. The Thai leukaemia cohort is the only record spanning three stages, and it is the
reason the adolescent window is non-empty at all. The in-utero window is empty for the
second consecutive issue.}

\vspace{2mm}
\noindent\includegraphics[width=\linewidth]{\FIGDIR/f4_heatmap.png}
\figcap{Subtopic $\times$ study-architecture cross-tabulation. The silo the 12 August
issue flagged is if anything sharper today: every \textit{Sensing} record sits in a
methods, instrument or modelling column and none carries a health endpoint, while every
health record sits in an observational column and none reports its own measurement.}

\vspace{2mm}
\noindent\includegraphics[width=\linewidth]{\FIGDIR/f5_forest.png}
\figcap{The four interval estimates carried today. \textbf{The two pairs are on
completely different exposure contrasts and must not be compared to each other}: Xu et
al.'s ORs are top-versus-bottom contrasts on a 3-year mean in a clinical population with
established intracranial stenosis --- a selected, high-risk sample in which effect sizes
above 4 are not implausible but are not population estimates --- while Tan et al.'s are
per-quantile mixture-index ORs from two different mixture estimators (WQS and QGC) fitted
to the same data, so they are not independent of each other either. Xu et al.'s
component-specific ORs for sulfate, organic matter, nitrate and ammonium are reported in
the abstract without intervals and are therefore not plotted.}

\clearpage

% ==================================================================== DIGEST
\section*{\sffamily\bfseries\color{Deep}Paper-level digest}
\vspace{-1.5mm}{\color{Amber}\rule{\linewidth}{1.1pt}}\vspace{1mm}

{\fontsize{7.4}{9}\selectfont\color{Slate}Ordered by cluster. Each entry gives the
methodological core and the finding that matters, not an abstract paraphrase.\par}

\band{Sensing, forecasting \& instrumentation\hfill 7 records}{Amber}

\paperentry{Benjamin et al.\ 2026}
{Limitations of environmental extrapolation in low-cost carbon monoxide and PM$_{2.5}$ sensors}
{Sensors (Basel)}
{Multi-season co-location of a Plantower PMS5003 and an Alphasense CO-B4 against Federal
Equivalence Method monitors, with calibration models deliberately trained on restricted
ranges of relative humidity, temperature and pollutant concentration and then scored
outside those ranges. \textbf{Within-envelope bias was minimal; crossing any envelope
wall introduced bias; crossing the upper concentration wall specifically increased random
error.} The design is the contribution --- most co-location papers report a single
in-sample $R^2$ and say nothing about the envelope. Limitations: one site, one sensor
model per pollutant, and no aerosol-composition contrast, so the concentration effect
cannot be separated from a possible composition shift at high concentration.}
{10.3390/s26154988}

\paperentry{Fawaz et al.\ 2026}
{Surface cleaning of SAW-based microparticle sensors in a cascade impactor via SAW-induced droplet displacement}
{Sensors (Basel)}
{Rayleigh waves generated on a lithium-niobate chip actuate a water droplet that sweeps
deposited particles out of the sensing zone of surface-acoustic-wave microparticle
sensors integrated into a cascade impactor. Optimal droplet size and required RF power
were determined; the droplet displaced both NaCl-solution residue and SiC particles
without degrading sensor performance. Surface saturation is the binding constraint on
long unattended deployment of mass-sensitive acoustic sensors, so a self-regenerating
cycle changes the maintenance model rather than the sensitivity. Laboratory
demonstration only --- no ambient aerosol, no duty-cycle endurance data, and no accounting
for what the droplet does to semi-volatile mass.}
{10.3390/s26154666}

\paperentry{Gao et al.\ 2026}
{Feasibility of an air sensor array for real-time detection and characterisation of VOCs}
{Atmos Meas Tech}
{Commercial VOC air sensors are broadband --- each returns one scalar reflecting an
ill-defined sensitivity envelope, which is why they have seen little research use. This
paper tests whether an \emph{array} spanning three fundamentally different sensor types
can recover speciation from the joint response. The measurand is VOCs, not PM, and it is
carried here for the method: the same argument applies directly to optical PM sensors
whose single scalar folds together size distribution, refractive index and hygroscopic
growth. Feasibility study --- laboratory characterisation, and the practical question of
whether the array's discrimination survives ambient interferent mixtures is not settled.}
{10.5194/amt-19-5337-2026}

\paperentry{Wu \& Qu 2026}
{Interpretable short-term PM$_{2.5}$ forecasting using meteorology--pollution coupling across Beijing stations}
{Environ Monit Assess}
{420{,}768 hourly observations from 12 Beijing stations, 2013--2017, under a strict
chronological train--validation--test split with Diebold--Mariano testing, feature-group
ablation and station-holdout transfer. At the 6\,h horizon CatBoost gave
\textbf{RMSE 51.24\,$\mu$g/m$^3$, MAE 30.87, $R^2=0.617$}, a 6.73\% RMSE reduction over
persistence; station-holdout LightGBM reached 53.11\,$\mu$g/m$^3$, 5.65\% below
persistence. Dominant predictors were recent PM history, wind speed and direction,
diurnal cyclicity and the temperature--dew-point gap. \textbf{The honest part is the
framing}: a 6.7\% gain over persistence is a modest operational return, performance
degrades during high-pollution episodes, and the authors state that the 2013--2017 regime
would require retraining for contemporary deployment.}
{10.1007/s10661-026-15799-9}

\paperentry{Adams et al.\ 2026}
{Sampling sub-pixel scale variability in total column NO$_2$ within a coastal urban environment}
{Environ Sci Atmos}
{Long-term Pandora spectrometer records of column-integrated NO$_2$ at Boston University
and Harvard University, used to characterise variability \emph{within} a single satellite
pixel. Directly relevant to any satellite--surface fusion product: sub-pixel
heterogeneity is an irreducible error term in the validation of column retrievals, and
two sites a few kilometres apart in the same coastal urban area are the cleanest way to
bound it. Two instruments in one metropolitan area is a small basis for a general claim
about pixel-scale variability, and NO$_2$ column heterogeneity is not a direct proxy for
PM heterogeneity.}
{10.1039/d6ea00033a}

\paperentry{Liu et al.\ 2026c}
{Polarization calibration of spaceborne lidar using dense cirrus-scattered solar background}
{Atmos Meas Tech}
{Solar background scattered by optically thick ice cloud provides a semi-continuous
daytime polarization calibration that complements onboard pseudo-depolarizers. At 1064\,nm
(CATS) molecular scattering is negligible; at shorter wavelengths it polarizes the
background and biases the calibration. The paper adds a vector-radiative-transfer
molecular-scattering correction and shows the daytime bias at CALIOP 532\,nm is
\textbf{$\sim$1\%}, removed by the correction to agree with onboard calibration; EarthCARE
ATLID at 355\,nm is treated by model calculation only. Carried at tier C: depolarization
ratio drives satellite aerosol typing, which in turn feeds AOD-to-PM$_{2.5}$ conversion,
but the link to surface concentration is two inferential steps away.}
{10.5194/amt-19-5325-2026}

\paperentry{Liu et al.\ 2026d}
{LRM-YOLO: lightweight YOLOv10n model for forest fire smoke detection in UAV images}
{Sensors (Basel)}
{YOLOv10n with a RepViTBlock feature-extraction module, a lightweight efficient
convolutional detection head, and MPDIoU bounding-box regression loss. Against baseline:
\textbf{$-$36.7\% parameters, $-$42.3\% GFLOPs, $+$1.3\% mAP50, $+$3.7\% mAP50-95},
recall $+$3.2\%, precision $+$3.6\%, on a purpose-built UAV-perspective dataset plus
Wildfire Smoke V1. Smoke detection in imagery is a plume-presence proxy, not a
concentration measurement, and the self-constructed dataset means the headline gains are
partly a property of the benchmark. Relevant only as an early-warning trigger upstream of
a real exposure measurement.}
{10.3390/s26154887}

\clearpage

\band{Exposure assessment \& modelling\hfill 2 records}{Clay}

\paperentry{Gould \& Jagarnath 2026}
{Satellite monitoring for air quality and public health research in Africa: a scoping review}
{Research Square (preprint)}
{JBI methodology, PRISMA-ScR reporting, five databases plus grey literature to May 2026,
for studies using satellite-derived air quality data against human health outcomes in
Africa. \textbf{Twenty-six records met inclusion} --- for an entire continent. Most
post-2015, most ecological time-series or cross-sectional, concentrated in West Africa,
with PM$_{2.5}$ the dominant exposure. This is the first record in this watch to put a
number on the African evidence gap that its own geography figure has shown as an empty
bar every day since 27 July. Preprint, not peer reviewed; a scoping review counts
literature and does not weigh it.}
{10.21203/rs.3.rs-10623736/v1}

\paperentry{da Costa Cardoso et al.\ 2026}
{Legacy Pb--Zn mine tailings as persistent sources of bioaccessible toxic elements in urban dust}
{Environ Monit Assess}
{18 urban dust samples and reference soils from Boquira, a former Pb--Zn district in
northeastern Brazil; microwave-assisted digestion and XRF for Pb, Zn, Cd, Cu, Mn, Fe, Al;
geoaccumulation index, contamination and enrichment factors, ecological risk index; oral
bioaccessibility by dilute-nitric-acid gastric simulation. \textbf{Pb bioaccessibility
$\sim$75\%}, Cu 54.8\%, Mn 28.1\%, Zn 27.0\%. Oral ingestion dominated over dermal;
non-carcinogenic HI exceeded 1 for children on Pb, and carcinogenic risk from Cd exceeded
limits for both children and adults. Eighteen samples in one district, a single gastric
extraction protocol without an intestinal phase, and no particle-size fractionation ---
which matters, because resuspension and inhalation-relevant dose are size-dependent and
this design cannot resolve them.}
{10.1007/s10661-026-15743-x}

\band{Neuro, cognition \& mental health\hfill 1 record}{Deep}

\paperentry{Xu et al.\ 2026}
{PM$_{2.5}$ components and cerebral hypoperfusion, with mediation by inflammatory biomarkers}
{Eur J Neurol}
{967 patients with symptomatic intracranial arterial stenosis at one Suzhou hospital;
hypoperfusion quantified from multimodal CT with MIStar; 3-year PM$_{2.5}$ and component
exposures from the TAP database; generalized additive models, restricted cubic splines,
weighted quantile sum for the mixture, and mediation analysis.
\textbf{PM$_{2.5}$ OR 4.344 (3.323--5.734); black carbon 4.396 (3.327--5.869)}, then
sulfate 4.359, organic matter 4.278, nitrate 3.552, ammonium 3.313 (no intervals given).
OM (44.2\%) and BC (42.2\%) carry 86\% of the mixture effect. Mediation by neutrophils
5.37\%, leukocytes 4.30\%, monocytes 2.73\% --- i.e.\ inflammation accounts for roughly a
tenth of the total effect, so the dominant pathway is unexplained. Single-centre clinical
sample already selected on stenosis; ORs above 4 in that population are not transportable
to general-population risk. Component exposures are modelled, not measured.}
{10.1111/ene.70716}

\band{Occupational \& indoor\hfill 1 record}{Clay}

\paperentry{Tan et al.\ 2026}
{PM$_{2.5}$ components, nutritional indices and leisure-time physical activity vs lower eGFR in subway workers}
{Nutrients}
{8{,}477 Wuhan subway workers; PM$_{2.5}$ and components from the Tracking Air Pollution
in China dataset; creatinine-based eGFR; GNRI and PNI nutritional indices; IPAQ-derived
activity. Mixture effects by WQS, quantile g-computation and BKMR.
\textbf{WQS OR 1.183 (1.106--1.266) and QGC OR 1.183 (1.083--1.291)} for lower eGFR;
BKMR overall effect on eGFR of \textbf{$-0.045$ ($-0.068$, $-0.022$)} for a 5-percentile
mixture increase against medians. Subgroups by nutrition and activity differed but
\textbf{no significant interaction} was found --- the authors say so, and the correct
reading is hypothesis generation, not effect modification. Cross-sectional, so reverse
causation is unaddressed; and residential exposure is a poor proxy for the occupational
exposure of people who work underground, which is the population's defining feature.}
{10.3390/nu18152580}

\band{Mechanistic toxicology\hfill 2 records}{Violet}

\paperentry{Gomes et al.\ 2026}
{Senolytics reverse PM$_{2.5}$-induced cell and vascular dysfunction}
{Cardiovasc Toxicol}
{C57BL/6J mice inhaling concentrated ambient particles versus filtered air showed
endothelial activation \emph{in vivo} --- increased leukocyte rolling on and adhesion to
intact vasculature --- alongside senescence of peripheral blood mononuclear cells and
endothelial progenitor cells. Dasatinib + quercetin reversed the senescence readouts
($\beta$-galactosidase activity, SASP gene expression), restored EPC tube-forming capacity
in vitro, and limited in vivo endothelial activation. Positions senescence upstream of
early atherogenesis rather than as a late consequence. Concentrated ambient particles
deliver a supra-ambient dose; the senolytic rescue demonstrates sufficiency of the
senescence pathway under that dose, not its necessity at ambient concentrations. No
exposure--response and no female animals reported.}
{10.1007/s12012-026-10171-2}

\paperentry{Zhang et al.\ 2026}
{Pleural anthracosis severity and the structure and function of subcarinal lymph nodes}
{J Thorac Dis}
{Thoracoscopic grading of pleural anthracosis into four grades by three thoracic surgeons
with ImageJ area quantification, against paired subcarinal lymph nodes stained with HE,
Masson, EVG and immunohistochemistry for Ki67, CD31, CD68 and D2-40. Black particulate
accumulation in node capsule and medulla rose with anthracosis grade, but every
quantitative immunostaining index \textbf{rose then fell} with advancing severity --- a
consistent non-monotonic pattern across proliferation, vascularity, macrophage presence
and lymphatic integrity. The non-monotonicity is the finding and is also the weakness:
with a graded surgical sample and no exposure history, an exhaustion interpretation and a
selection interpretation are equally consistent with the data.}
{10.21037/jtd-2026-1274}

\clearpage

\band{Cardiovascular \& metabolic\hfill 2 records}{Teal}

\paperentry{Chen et al.\ 2026}
{Compound cold wave--PM$_{2.5}$ exposure and cardiovascular emergency ambulance calls}
{Research Square (preprint)}
{82{,}714 cardiovascular emergency ambulance call records from Shijiazhuang, 1 January
2014 to 31 December 2023, with three cold-wave categories defined on temperature
percentiles crossed with four PM$_{2.5}$ pollution levels to construct compound exposure
events; distributed lag non-linear model with quasi-Poisson regression for lagged and
interactive effects. The compound-exposure framing is the contribution --- most
temperature $\times$ PM work tests a multiplicative interaction term rather than
constructing joint exposure states. Preprint, not peer reviewed; ambulance calls are a
severity-filtered and access-dependent outcome; and a 3$\times$4 exposure grid over ten
years leaves the extreme cells thin.}
{10.21203/rs.3.rs-10547139/v1}

\paperentry{Gaggini \& Vassalle 2026}
{Gene--air pollution interaction in cardiovascular disease}
{Int J Mol Sci}
{Narrative PubMed review of combined genetic-susceptibility and air-pollution effects on
cardiovascular risk, concentrating on detoxification and inflammation genes. Argues the
area could refine cardiovascular risk prediction beyond traditional factors. Narrative,
not systematic; gene--environment interaction in this literature is dominated by
underpowered candidate-gene studies, which the review acknowledges as ``lights and
shadows'' but does not quantify. Useful only as an entry point to the literature.}
{10.3390/ijms27156818}

\band{Other clinical endpoints\hfill 1 record}{Slate}

\paperentry{Traisathit et al.\ 2026}
{Air pollution and mortality in acute lymphoblastic leukaemia across age groups, northern Thailand}
{PLoS One}
{604 ALL patients diagnosed 1999--2020 in upper northern Thailand --- a region with severe
seasonal biomass-burning haze --- stratified into paediatric ($<$15, n=344), adolescent and
young adult (15--39, n=152) and adult ($\geq$40, n=108). Five-year survival 64.5\%,
19.6\% and 6.7\% respectively; adjusted HR for death 3.76 (AYA) and 8.10 (adult) versus
paediatric; diagnosis before 2010 aHR 1.38. \textbf{PM$_{2.5}$ $\geq$50\,$\mu$g/m$^3$ was
associated with increased mortality in paediatric patients.} The age-group HRs are
prognosis, not pollution. The pollution claim rests on comparing mean concentrations
between those who died and those who did not, with a threshold analysis in one stratum;
no exposure window, no treatment-era adjustment beyond the pre/post-2010 term, and
ambient station concentrations as individual exposure. Signal, not estimate.}
{10.1371/journal.pone.0355856}

\band{Burden, policy \& mitigation\hfill 1 record}{Coral}

\paperentry{He et al.\ 2026}
{Premature ischaemic stroke burden in adults aged 50--69 in China, India and the United States}
{Neuroepidemiology}
{GBD 2023 DALYs, YLLs, YLDs and population attributable fractions for the 50--69 stratum
in three countries at different stages of epidemiological transition, with ARIMA
projection to 2053. Trajectories diverge: China declining, India flat at a high level,
the US on a high plateau; the US has shifted to a YLD-dominant (disability) composition
while India remains YLL-dominant (acute mortality); China shows a near 2:1 male:female
ratio driven by male smoking. \textbf{India is characterised as facing dual pressure from
persistent PM$_{2.5}$ and rising metabolic risk.} The authors are explicit that widening
intervals make the 2053 horizon scenario-based rather than forecast; a 34-year ARIMA
extrapolation from modelled burden estimates is presentation, not evidence. Carried for
the PAF decomposition only.}
{10.1159/ned/adwag005}

\clearpage

% ==================================================================== TABLE
\section*{\sffamily\bfseries\color{Deep}Corpus register}
\vspace{-1.5mm}{\color{Amber}\rule{\linewidth}{1.1pt}}\vspace{2.5mm}

{\fontsize{6.7}{8.4}\selectfont
\setlength{\tabcolsep}{2.4pt}
\begin{longtable}{@{}>{\RaggedRight}p{25mm}>{\RaggedRight}p{41mm}>{\RaggedRight}p{27mm}%
>{\RaggedRight}p{23mm}>{\RaggedRight}p{18mm}>{\RaggedRight}p{15mm}@{}}
\toprule
\textsf{\textbf{First author}} & \textsf{\textbf{Focus}} & \textsf{\textbf{Journal}} &
\textsf{\textbf{Design}} & \textsf{\textbf{Metric}} & \textsf{\textbf{Region}}\\
\midrule
\endfirsthead
\toprule
\textsf{\textbf{First author}} & \textsf{\textbf{Focus}} & \textsf{\textbf{Journal}} &
\textsf{\textbf{Design}} & \textsf{\textbf{Metric}} & \textsf{\textbf{Region}}\\
\midrule
\endhead
\multicolumn{6}{@{}l@{}}{\cellcolor{Mist}\textsf{\textbf{\textcolor{Deep}{Neuro / mental health}}}}\\[0.6mm]
\href{https://doi.org/10.1001/jamanetworkopen.2026.25015}{Dales} & Wildfire-sourced PM$_{2.5}$ and ED visits for migraine and other primary headache & \textit{JAMA Netw Open} & Case-crossover & PM$_{2.5}$ (wildfire) & Canada\\
\href{https://doi.org/10.1093/ageing/afag215}{Glover} & Umbrella meta-review of urban environmental exposures and cognitive health & \textit{Age Ageing} & Umbrella review / meta-analysis & PM$_{2.5}$ & Global\\
\href{https://doi.org/10.1016/j.nsa.2026.107019}{Goossens} & In utero air pollution exposure and pre-/postnatal brain development & \textit{Neurosci Appl} & Narrative review & PM$_{2.5}$ / PM$_{10}$ / PM$_{\mathrm{coarse}}$ & Global\\
\href{https://doi.org/10.1016/j.isci.2026.116617}{Ma} & PM$_{2.5}$-induced neurotoxicity and depression: exposure risks and mechanisms & \textit{iScience} & Narrative review & PM$_{2.5}$ + components & Global\\
\href{https://doi.org/10.3760/cma.j.cn112150-20250804-00752}{Pan} & Maternal PM exposure and offspring neurodevelopment: effects and mechanisms & \textit{Zhonghua Yu Fang Yi Xue Za Zhi} & Narrative review & PM (bulk) & China\\
\addlinespace[0.9mm]
\multicolumn{6}{@{}l@{}}{\cellcolor{Mist}\textsf{\textbf{\textcolor{Teal}{Cardiovascular \& metabolic}}}}\\[0.6mm]
\href{https://doi.org/10.3760/cma.j.cn112338-20251120-00843}{Huang} & PM$_{2.5}$ and dyslipidaemia in rural Chinese children; dietary effect modification & \textit{Zhonghua Liu Xing Bing Xue Za Zhi} & Cross-sectional (multistage) & PM$_{2.5}$ & China\\
\href{https://doi.org/10.32394/pe/218966}{Kazmierska} & Air pollution and cardiovascular disease: review of current data & \textit{Przegl Epidemiol} & Narrative review & PM$_{2.5}$ / PM$_{10}$ & Global\\
\href{https://doi.org/10.1007/s00125-026-06814-2}{Xi} & Air pollution, built environment and temperature vs. prediabetes phenotypes & \textit{Diabetologia} & Prospective cohort & PM$_{2.5}$ & Germany\\
\href{https://doi.org/10.1161/JAHA.125.046917}{Yang} & Short-term PM$_{2.5}$ and hospitalisation/ED use in atrial fibrillation & \textit{J Am Heart Assoc} & Case-crossover & PM$_{2.5}$ (wildfire-influenced) & USA\\
\addlinespace[0.9mm]
\multicolumn{6}{@{}l@{}}{\cellcolor{Mist}\textsf{\textbf{\textcolor{Sky}{Reproductive \& developmental}}}}\\[0.6mm]
\href{https://doi.org/10.1016/j.jes.2025.10.043}{Chen} & Prenatal air pollution, fetal growth trajectory phenotypes and child neurodevelopment & \textit{J Environ Sci (China)} & Prospective cohort & PM$_{2.5}$ / PM$_{10}$ & China\\
\href{https://doi.org/10.1016/j.ecoenv.2026.120566}{Huang} & Joint prenatal bisphenol + air-pollutant mixture exposure and preschool IQ & \textit{Ecotoxicol Environ Saf} & Prospective cohort (mixtures) & PM$_{2.5}$ / PM$_{10}$ & China\\
\href{https://doi.org/10.1007/s40572-026-00552-8}{Peterson} & Ambient air pollution, PAHs and pubertal development: critical review & \textit{Curr Environ Health Rep} & Critical review & PM$_{2.5}$ components & Global\\
\href{https://doi.org/10.1016/j.envres.2026.125313}{Xie} & Long-term PM$_{2.5}$ and timing of pubertal development in boys and girls & \textit{Environ Res} & Prospective cohort & PM$_{2.5}$ & China\\
\addlinespace[0.9mm]
\multicolumn{6}{@{}l@{}}{\cellcolor{Mist}\textsf{\textbf{\textcolor{Sage}{Respiratory \& allergic}}}}\\[0.6mm]
\href{https://doi.org/10.1007/s00484-026-03289-z}{Bayar Muluk} & Climate change and otolaryngology: airway disease impacts and emission reduction & \textit{Int J Biometeorol} & Narrative review & PM$_{2.5}$ & Global\\
\href{https://doi.org/10.4103/mgr.MEDGASRES-D-25-00164}{Liao} & Short-term air pollution, respiratory symptoms and lung function in asthma & \textit{Med Gas Res} & Panel study & PM$_{2.5}$ / PM$_{10}$ & China\\
\href{https://doi.org/10.3390/epidemiologia7040097}{Sanchez de Prada} & Ambient PM and asthma-coded admissions in children 0-3 y, Spain & \textit{Epidemiologia} & Ecological panel & PM$_{2.5}$ / PM$_{10}$ & Spain\\
\href{https://doi.org/10.1016/j.jhazmat.2026.143076}{Yang} & Allergenic vs. toxic risk of airborne seed hairs compared with ambient PM$_{2.5}$ & \textit{J Hazard Mater} & Field measurement & PM$_{2.5}$ + bioaerosol & China\\
\addlinespace[0.9mm]
\multicolumn{6}{@{}l@{}}{\cellcolor{Mist}\textsf{\textbf{\textcolor{Violet}{Mechanistic toxicology}}}}\\[0.6mm]
\href{https://doi.org/10.1016/j.jes.2025.12.074}{Ding} & PM$_{2.5}$ promotes vascular calcification via miR-27b-5p/HES1 signalling & \textit{J Environ Sci (China)} & In vitro + organoid + murine & PM$_{2.5}$ & China\\
\href{https://doi.org/10.1007/s00360-026-01698-5}{Maraschi} & Osmoregulatory disruption and metal bioaccumulation from settleable atmospheric PM & \textit{J Comp Physiol B} & Controlled exposure (ecotox) & Settleable APM & Brazil\\
\href{https://doi.org/10.1016/j.jes.2025.11.018}{Wu} & Macrophage NLRP3 inflammasome in real-ambient PM-induced hepatic glucose dysregulation & \textit{J Environ Sci (China)} & Real-ambient murine + co-culture & Real-ambient PM & China\\
\addlinespace[0.9mm]
\multicolumn{6}{@{}l@{}}{\cellcolor{Mist}\textsf{\textbf{\textcolor{Clay}{Occupational \& indoor}}}}\\[0.6mm]
\href{https://doi.org/10.3389/fmedt.2026.1796828}{Chen} & Nested suction instrument to cut intraoperative surgical-smoke PM$_{2.5}$ & \textit{Front Med Technol} & Bench / device experiment & PM$_{2.5}$ (surgical smoke) & China\\
\href{https://doi.org/10.1093/annweh/wxag055}{Cole} & Determinants of respirable crystalline silica exposure in Australian tunnelling & \textit{Ann Work Expo Health} & Observational exposure study & Respirable dust / RCS & Australia\\
\href{https://doi.org/10.1016/j.jhazmat.2026.143081}{Ding} & Stage-resolved hair Zn-isotope dynamics as a firefighting combustion-exposure tracer & \textit{J Hazard Mater} & Repeated-measures biomarker & Combustion-derived PM & China\\
\href{https://doi.org/10.3390/jox16040119}{Ramsden} & Xenobiotic hazards in aircraft cabin air & \textit{J Xenobiot} & Review + risk economics & Ultrafine particles & Global\\
\addlinespace[0.9mm]
\multicolumn{6}{@{}l@{}}{\cellcolor{Mist}\textsf{\textbf{\textcolor{Amber}{Sensing, forecasting \& instrumentation}}}}\\[0.6mm]
\href{https://doi.org/10.1016/j.jenvman.2026.130515}{Shi} & Adaptive spatiotemporal graph transformer for multi-site multi-step PM$_{2.5}$ forecasting & \textit{J Environ Manage} & Deep-learning model & PM$_{2.5}$ & China / India\\
\href{https://doi.org/10.3390/bios16070380}{Tsuchiya} & Seasonal and indoor-outdoor characteristics of fluorescent aerosol particles in offices & \textit{Biosensors} & Real-time sensor campaign & Bioaerosol / coarse & Japan\\
\addlinespace[0.9mm]
\multicolumn{6}{@{}l@{}}{\cellcolor{Mist}\textsf{\textbf{\textcolor{Amber}{Exposure assessment \& modelling}}}}\\[0.6mm]
\href{https://doi.org/10.1029/2025GH001621}{Solorzano-Araque} & Multiscale exposure modelling: LOTOS-EUROS + mobility surveys + particle toxicity & \textit{GeoHealth} & CTM + mobility integration & PM$_{2.5}$ / PM$_{10}$ & Colombia\\
\addlinespace[0.9mm]
\multicolumn{6}{@{}l@{}}{\cellcolor{Mist}\textsf{\textbf{\textcolor{Coral}{Burden, policy \& mitigation}}}}\\[0.6mm]
\href{https://doi.org/10.1002/hsr2.72886}{Karamzad} & Ischaemic heart disease burden in MENA 1990-2021; PM attributable fraction 28.2\% & \textit{Health Sci Rep} & GBD secondary analysis & Ambient PM & MENA\\
\href{https://doi.org/10.1371/journal.pone.0347989}{Rees} & HEPA filtration to reduce respiratory infection in care homes & \textit{PLoS One} & RCT process evaluation & Filtered PM & UK\\
\href{https://doi.org/10.1016/j.jes.2026.01.009}{Shan} & Policy-driven emission reductions and post-pandemic rebound, high-resolution inventories & \textit{J Environ Sci (China)} & Emission inventory & PM emissions & China\\
\addlinespace[0.9mm]
\multicolumn{6}{@{}l@{}}{\cellcolor{Mist}\textsf{\textbf{\textcolor{Slate}{Other clinical endpoints}}}}\\[0.6mm]
\href{https://doi.org/10.3390/dj14070440}{Lupita} & PM$_{2.5}$, preventive-care capacity and caries/edentulism burden across the EU & \textit{Dent J} & Ecological panel & PM$_{2.5}$ & EU-27\\
\href{https://doi.org/10.1016/j.jdin.2026.06.008}{Park} & Long-term PM exposure and risk of androgenetic alopecia & \textit{JAAD Int} & Nationwide cohort & PM$_{2.5}$ / PM$_{10}$ & South Korea\\
\href{https://doi.org/10.3171/2026.3.PEDS2615}{Sun} & Environmental exposure and area deprivation vs. paediatric hydrocephalus outcomes & \textit{J Neurosurg Pediatr} & Registry cohort + geospatial & PM$_{2.5}$ (satellite-derived) & USA\\
\addlinespace[0.9mm]
\bottomrule
\end{longtable}}


\vspace{4mm}
\begin{keybox}
{\sffamily\bfseries\footnotesize\color{Deep}Provenance, screening and known gaps}\\[1.2mm]
{\fontsize{7.2}{8.8}\selectfont
\textbf{Window.} \texttt{2026-08-13 -> 2026-08-13}, contiguous with the 12 August issue;
no gap. PubMed is windowed on \texttt{[Date - Entry]}, the Crossref-by-ISSN leg on
Crossref \texttt{created}, and the Europe PMC leg on \texttt{CREATION\_DATE}. Three
different clocks, not treated as interchangeable.\\[0.8mm]
\textbf{Per-leg yield.} PubMed connector, two queries (health axis 10, instrumentation
axis 4, one overlapping): 13 distinct, \textbf{9 carried}. Local harvester:
\texttt{pubmed\_health} 12 and \texttt{pubmed\_sensing} 4 surfaced \textbf{4 further
carried records} the connector queries missed --- the Sensors papers on LCS extrapolation,
SAW sensor cleaning and UAV smoke detection --- which is the clearest case yet for keeping
both legs. Europe PMC 2 (\textbf{2 carried}, both Research Square preprints). Crossref by
ISSN 8 (\textbf{3 carried}; three of the eight were RSC front matter --- contents list,
front cover, back cover). arXiv returned 145\,kB with nothing in scope on or after
6 August. \textbf{OpenAlex HTTP 429, 11th consecutive.} Consensus was not re-queried this
issue: the 12 August sweep on the calibration axis returned only out-of-window results
and this is the sixth consecutive issue with zero Consensus carry. Scholar Gateway's
corpus ends May 2026 and cannot serve an August window. ClinicalTrials.gov: no PM or
air-pollution registration updated in the window; \texttt{trials.json} unchanged.\\[0.8mm]
\textbf{Rejections (12).} Off-topic: broiler-diet life cycle assessment; bronchoscopic
lavage case series (its ``particulate matter'' is aspirated gastric content); tourism and
health satisfaction in older Chinese adults (air quality is one covariate among many);
indirect soil NO$_x$ emission budget; ground-based hyperspectral retrieval of
\emph{riverine} particulate organic carbon; zinc-pyrithione ultrafine particle synthesis
in a microchannel reactor; Bayesian latent-component analysis of Volta Lake
\emph{sediment} PAHs; UK methane isotope source attribution; sodium-layer lidar
dynamics; and three RSC journal front-matter items. All logged with reasons to
\texttt{state/rejected.jsonl}.\\[0.8mm]
\textbf{Two records are preprints} (Research Square) and are labelled as such in the
register. They are carried because both fill a documented gap in this watch --- African
satellite monitoring and compound cold-wave $\times$ PM exposure --- and neither is used to
support a numeric claim in the analytics.\\[0.8mm]
\textbf{One DOI warning, accepted.} \texttt{10.1159/ned/adwag005} (He et al.) resolves at
doi.org but has no Crossref record; the Karger \texttt{ned/} prefix pattern is new. The
link was confirmed to resolve before shipping.\\[0.8mm]
\textbf{Labels.} Subtopic, design, particle-metric and geography labels are drawn only
from vocabularies already mapped in \texttt{plots.py} and \texttt{mktable.py}; no new leaf
label was written this issue. Assignment is single-label by dominant contribution and is
a judgement, not a controlled vocabulary.\\[0.8mm]
Content derived from PubMed, Europe PMC and Crossref metadata; abstracts remain the
copyright of their respective publishers. All DOI links resolve to the publisher of
record.\par}
\end{keybox}

\end{document}
